%% -*- coding:utf-8 -*-
\documentclass[output=book
%                ,series=tbls,
%		,modfonts
%		,nonflat
%		,multiauthors
%	        ,collection
%	        ,collectionchapter
%	        ,collectiontoclongg
 	        ,biblatex  
                ,decapbib 
                ,babelshorthands
%                ,showindex
                ,newtxmath
%                ,colorlinks, citecolor=brown % for drafts
%                ,draftmode
% 	        ,coverus
                ,uniformtopskip
%                ,tblseight % to get the fonts for the Chinese title of GT
		  ]{langscibook} 


\usepackage{ifthen}
\provideboolean{draft}
\setboolean{draft}{false}

\newtoggle{draft}\togglefalse{draft}
%\newtoggle{draft}\toggletrue{draft}
%\newtoggle{finished}\togglefalse{finished}

% uncomment if biblatex causes problems
%\hypersetup{draft} 
% check for NoHyper in the final version

% load memoize before german.sty

%\usepackage{nomemoize} 
\usepackage{memoize} % The figures have been created with an earlier version of texlive. Recreating
                     % them changes the layout of the pages, so we have to rely on them. Unless ... 16.03.2023
%\memoizeset{readonly}
% see below for further settings

\input{localmetadata.tex}
\input{localpackages.tex}
\input{localcommands.tex}
\input{localhyphenation.tex}


\input{locallangscifixes.tex}


% use germanic.tex to play with this, so that main.tex is always the correct file for final compilations
%\hypersetup{draft} % can be used to find errors with biblatex references linked to two different
                   % pages 10.07.2025

\mmzset{
  prefix={germanic.memo.dir/},
%  register=\todo{O{}+m},
%  prevent=\todo,
}


% \renewcommand{\itd}[1]{}
% \renewcommand{\itdobl}[1]{}
% \renewcommand{\itdopt}[1]{}

% \renewcommand{\addpages}{}
% \renewcommand{\addsource}{}
% \renewcommand{\addglosses}{}


%\bibliography{germanic}
\bibliography{bib-abbr,biblio-xdata-en,biblio}


\input germanic-include

\end{document}

%%% Local Variables:
%%% TeX-command-extra-options: "-shell-escape"
%%% End:
